\chapter{Prowadzenie projektu i dokumentacja}\label{ch:projekt}

Ten rozdział opisuje, jak prowadzić projekt misji i dokumentację eksperymentu. Praktyki pochodzą z programu REXUS/BEXUS (ESA, DLR, SNSA), w którym zespół Stardust uczestniczył~\cite{praca2022,sed}. Program prowadzi zespoły studenckie przez pełny cykl misji kosmicznej. Uproszczona wersja tych praktyk działa dobrze w misjach własnych.

\begin{sekcja}{Cykl projektu}
Podziel projekt na etapy. Na końcu każdego etapu zrób przegląd: zespół przedstawia stan projektu osobom spoza zespołu (absolwentom projektu, opiekunowi koła, zespołowi badawczemu). Przegląd kończy się listą działań z osobami odpowiedzialnymi i terminami.

\renewcommand{\arraystretch}{1.3}
\small
\begin{tabularx}{\linewidth}{@{}P{30mm}LL@{}}
\toprule
\textbf{Etap} & \textbf{Pytanie przeglądu} & \textbf{Wynik}\\\midrule
Koncepcja & Co chcemy zbadać i dlaczego? Czy to jest możliwe? & Cel, wymagania, wstępny budżet masy.\\
Projekt wstępny (PDR) & Czy koncepcja spełnia wymagania? & Schemat systemu, wybór komponentów, plan testów.\\
Projekt krytyczny (CDR) & Czy projekt jest gotowy do budowy? & Kompletny projekt, budżety, rejestr ryzyk.\\
Integracja i testy & Czy zbudowany sprzęt działa? & Wyniki testów (F5), poprawki.\\
Gotowość do lotu (FRR) & Czy możemy bezpiecznie polecieć? & Decyzja „leci / nie leci” (F6).\\
Po misji & Co działało? Co poprawić? & Raport, wnioski (F14).\\
\bottomrule
\end{tabularx}
\normalsize

Dla prostej misji szkoleniowej wystarczą dwa przeglądy: przed budową (CDR) i przed lotem (FRR).
\end{sekcja}

\begin{sekcja}{Dokumentacja eksperymentu}
Każdy ładunek i każda zmiana platformy ma dokument eksperymentu. Wzór: dokument SED (\emph{Student Experiment Documentation}) z programu REXUS/BEXUS. Przykład: dokument SED zespołu Stardust~\cite{sed}. Opis platformy BEXUS i wymagań wobec eksperymentów: BEXUS User Manual~\cite{bexusmanual}. Szablony przeglądów PDR, CDR i FRR dla studenckich ładunków balonowych udostępnia kurs LaACES~\cite{laaces}.

Minimalna struktura dokumentu eksperymentu:
\begin{enumerate}
  \item \textbf{Wprowadzenie} -- cel naukowy albo techniczny, zespół i role.
  \item \textbf{Wymagania} -- co eksperyment musi robić, jakie ma ograniczenia (masa, wymiary, zasilanie, temperatura).
  \item \textbf{Projekt} -- mechanika, elektronika, oprogramowanie (np. maszyna stanów).
  \item \textbf{Budżety} -- masa, energia, dane.
  \item \textbf{Rejestr ryzyk} -- co może się nie udać, prawdopodobieństwo, skutek, środek zaradczy.
  \item \textbf{Plan testów i wyniki} -- każde wymaganie ma test, który je weryfikuje.
  \item \textbf{Procedury} -- montaż, przygotowanie na stanowisku, formularze.
  \item \textbf{Wyniki i wnioski} -- po misji.
\end{enumerate}

Pisz dokument od początku projektu, nie po misji. Numeruj wersje. Zapisuj, co się zmieniło.
\end{sekcja}

\begin{sekcja}{Dokumentacja i kod}
Dokumentacja może być w repozytorium kodu albo w osobnym pliku. Ważne są dwa powiązania:
\begin{itemize}
  \item \textbf{dokumentacja w wydaniu kodu}: do wydania (\emph{release}) kodu lotnego dołącz końcową wersję dokumentacji i raport z misji,
  \item \textbf{wersja kodu w dokumentacji}: w karcie misji i w raporcie zapisz numer wydania albo commita kodu, który poleciał.
\end{itemize}
Ten podręcznik powstaje automatycznie z plików \LaTeX{} w repozytorium \www{github.com/simle-stardust/Manual}. Każde wydanie (tag \texttt{v*}) ma dołączony PDF.
\end{sekcja}

\begin{sekcja}{Budżety}
\textbf{Budżet masy:} tabela wszystkich elementów z masą zmierzoną. Zapas co najmniej \SI{10}{\percent}.

\textbf{Budżet energii:} dla każdego urządzenia: napięcie, prąd w pracy, prąd szczytowy, czas pracy. Suma energii $\times$ zapas $\leq$ pojemność baterii w niskiej temperaturze.

\textbf{Budżet danych:} dane na sekundę $\times$ czas lotu $\leq$ pojemność pamięci. Przepustowość łącza radiowego $\geq$ dane telemetrii.

\textbf{Budżet finansowy:} koszty jednorazowe (powłoka, hel, materiały, baterie, subskrypcje), koszty sprzętu wielokrotnego, transport, ubezpieczenie, zapas na powtórkę lotu. Budżet finansowy jest podstawą wniosku o dofinansowanie i oferty dla klienta (np. zespołu badawczego albo firmy, której ładunek leci). Zapisuj rzeczywiste koszty każdej misji -- następny budżet będzie dokładniejszy.
\end{sekcja}

\begin{sekcja}{Rejestr ryzyk}
\renewcommand{\arraystretch}{1.3}
\small
\begin{tabularx}{\linewidth}{@{}P{40mm}P{40mm}L@{}}
\toprule
\textbf{Ryzyko} & \textbf{Skutek} & \textbf{Środek zaradczy}\\\midrule
Odcinacz nie działa & Długi lot, lądowanie daleko & Testy w warunkach lotu, zapas linki, drugi odcinacz innego typu\\
Brak telemetrii & Utrata ładunku & Dwa systemy śledzenia, SPOT, znaczniki Bluetooth\\
Za mało gazu & Opóźniony start & Zapas gazu, pomiar siły nośnej\\
Rozładowana bateria & Brak danych i lokalizacji & Budżet energii, test w zimnie, nowe baterie\\
Lądowanie w wodzie lub za granicą & Utrata ładunku, incydent & Predykcja, kryteria braku startu, odcinacz\\
Restart OBC w locie & Utrata stanu misji & Stan w pamięci nieulotnej\\
Uszkodzenie przy lądowaniu & Utrata danych & Wkładka, mocowanie kart pamięci\\
\bottomrule
\end{tabularx}
\normalsize
\end{sekcja}

\begin{sekcja}{Nazwa, kryptonim i patch misji}\label{sec:patch}
Nadaj każdej misji nazwę i kryptonim, na przykład Antares (2018) albo Stormdust (2023). Używaj ich w dokumentach, nazwach plików, formularzach i w mediach -- wtedy misje w archiwum się nie mylą. Zaprojektuj też patch misji: naszywkę z nazwą, rokiem i symbolem celu misji (przykład: patch z 2023 roku na okładce). Patch łączy zespół, oznacza zdjęcia, prezentacje i koszulki oraz jest gotową grafiką do wpisów i podziękowań dla sponsorów. Zapisz nazwę, kryptonim i plik patcha w karcie misji (formularz F0).
\end{sekcja}

\begin{sekcja}{Zdjęcia i dokumentacja wizualna}
Zdjęcia to ważna część procedur. Służą do:
\begin{itemize}
  \item dokumentacji projektu -- jak sprzęt wyglądał w dniu startu,
  \item analizy po locie -- zdjęcia pokazują, co poszło nie tak (np. ułożenie linek, mocowanie anteny),
  \item promocji -- zdjęcia przyciągają sponsorów i nowych członków.
\end{itemize}

Wydziel osobną rolę: \textbf{fotograf / dokumentalista} (formularz F1). Fotograf nie ma innych zadań przy starcie. Fotograf fotografuje:
\begin{itemize}
  \item każdą gondolę otwartą i zamkniętą, przed zamknięciem,
  \item stanowisko, tankowanie, zamknięcie szyjki, zestaw przed startem,
  \item tankowanie i start -- w całości filmem. Przy anomalii film pozwala ustalić jej przyczynę~\cite[rozdz.~11]{nasahb},
  \item miejsce lądowania przed dotknięciem sprzętu.
\end{itemize}
\end{sekcja}

\begin{sekcja}{Raport z misji}
Napisz raport w ciągu tygodnia po misji. Wzór: raport z misji testowej 2020~\cite{raport2020}.
\begin{enumerate}
  \item Zestaw lotny: elementy, masy, wersje kodu.
  \item Oś czasu: NOTAM, start, pęknięcie albo odcięcie, lądowanie, odzyskanie.
  \item Dane z lotu: trasa, profil wysokości, czujniki (pliki w załączniku).
  \item Problemy: co nie działało i dlaczego.
  \item Wnioski: co zmienić w następnej misji.
  \item Plan zmian: co, kto, do kiedy~\cite[rozdz.~11]{nasahb}.
\end{enumerate}
Dopisz wnioski do rozdziału~\ref{ch:historia}.
\end{sekcja}

\begin{sekcja}{Przekazywanie wiedzy}
Zespoły studenckie zmieniają się co kilka lat. Wiedza ginie, gdy jest tylko w głowach.
\begin{itemize}
  \item Każda nowa grupa zaczyna od lotu szkoleniowego z prostym ładunkiem.
  \item Doświadczona osoba prowadzi pierwszą misję. Następną misję nowy zespół prowadzi samodzielnie.
  \item Każdy podzespół ma osobę odpowiedzialną i opisaną procedurę.
  \item Po każdej misji aktualizuj tę dokumentację.
  \item Publikuj wyniki (konferencje, artykuły, blog). Publikacje zostają po zespole.
\end{itemize}
\end{sekcja}

\begin{sekcja}{Dalsza lektura}
\begin{itemize}
  \item \textbf{NASA Scientific Ballooning Handbook} (wydanie na 50-lecie, 2025) -- kompendium wiedzy o balonach naukowych: teoria lotu, gazy nośne, operacje, spadochrony, łączność. Na początek przeczytaj rozdział~2 o balonach w edukacji (programy studenckie, wnioski zespołów BOREALIS i NEBP)~\cite{nasahb}.
  \item \textbf{Projekt Loon} (X, the moonshot factory) -- doświadczenia z tysięcy lotów balonów stratosferycznych: konstrukcja powłok, nawigacja z wiatrem, operacje. Kolekcja materiałów i podsumowanie „The Loon Library”~\cite{loon}.
  \item \textbf{BEXUS User Manual} -- opis platformy BEXUS i wymagań wobec eksperymentów~\cite{bexusmanual}.
  \item \textbf{Kursy balonowe dla studentów}: LaACES (elektronika, programowanie, zarządzanie projektem, szablony dokumentów)~\cite{laaces} oraz kursy projektu NEBP (planowanie misji i logistyka startu)~\cite{nebp}.
  \item \textbf{Akademia Stratosferycznie} -- bezpłatne kursy o misjach stratosferycznych po polsku~\cite{akademia}.
  \item Artykuły zespołu Stardust z konferencji IAC i ESA PAC~\cite{iac2023,pac2024}.
\end{itemize}
\end{sekcja}
